\chapter{L'interface et la démonstration}
\label{ch:demo}

\section{L'interface}

Une seule page Streamlit, trois zones. La conversation à gauche ; à droite, le profil du client, le
panneau \og pourquoi \fg{} et le panneau \og frugalité \fg. Le jury doit pouvoir lire la décision et
son coût sans qu'on ait à les commenter.

\begin{figure}[H]
\centering
\begin{adjustbox}{max width=\textwidth}
\begin{tikzpicture}[x=1cm,y=1cm,font=\scriptsize]
  % Cadre
  \draw[darkgray,rounded corners=2pt,fill=white] (0,0) rectangle (15.4,-9.4);
  \fill[mainblue] (0,0) rectangle (15.4,-0.7);
  \node[white,anchor=west,font=\small\bfseries] at (0.25,-0.35) {Assistant bancaire --- prototype};
  \node[white,anchor=east] at (15.2,-0.35)
    {[ Client : Léa (C07) $\blacktriangledown$ ] \quad [ Mots simples ] \quad [ A+ ]};
  \fill[ext!70] (0,-0.7) rectangle (15.4,-1.15);
  \node[anchor=west] at (0.25,-0.92) {Vous échangez avec un assistant automatique. Il oriente ;
    un conseiller décide.};
  % Conversation
  \draw[darkgray!60] (0.25,-1.4) rectangle (9.2,-9.15);
  \node[anchor=north west,font=\scriptsize\bfseries,text=darkgray] at (0.35,-1.45) {Conversation};
  \node[acteur,anchor=east,text width=4.6cm,align=left] at (9.0,-2.25)
    {Je veux acheter un appartement};
  \node[etage,anchor=west,text width=6.2cm,align=left] at (0.45,-3.35)
    {C'est un beau projet. J'ai noté vos revenus de 2\,900\eur{} nets par mois. Quel est le prix
    du bien que vous visez, et quelle part de vos 25\,000\eur{} d'épargne souhaitez-vous y
    consacrer ?};
  \node[acteur,anchor=east,text width=4.6cm,align=left] at (9.0,-4.55)
    {Autour de 250\,000, et je peux mettre les 25\,000};
  \node[etage,anchor=west,text width=6.2cm,align=left] at (0.45,-5.95)
    {Avec ce prix, la mensualité serait d'environ 1\,220\eur{}, soit 42\,\% de vos revenus : au-delà
    des 35\,\% retenus pour un crédit immobilier. Avec votre apport, un bien autour de
    212\,000\eur{} resterait dans cette limite. \underline{Taux d'effort} \dots};
  \node[module,anchor=west,text width=6.2cm,align=left] at (0.45,-7.45)
    {\textbf{Action proposée} : prendre rendez-vous avec un conseiller \quad [ Oui ] [ Non ]};
  \draw[darkgray!60] (0.45,-8.45) rectangle (9.0,-8.95);
  \node[anchor=west,text=darkgray] at (0.55,-8.7) {Écrivez votre message\dots};
  % Profil
  \draw[darkgray!60] (9.45,-1.4) rectangle (15.15,-3.45);
  \node[anchor=north west,font=\scriptsize\bfseries,text=darkgray] at (9.55,-1.45) {Profil};
  \node[anchor=north west,text width=5.4cm,align=left] at (9.55,-1.8)
    {Léa, 31 ans, cadre\\
     Aucun crédit en cours\\
     Revenus : \textit{2\,900\eur{} nets par mois}\\
     Épargne : \textit{25\,000\eur{}}\quad{\tiny\textit{(italique : simulé)}}};
  % Pourquoi
  \draw[darkgray!60] (9.45,-3.6) rectangle (15.15,-6.95);
  \node[anchor=north west,font=\scriptsize\bfseries,text=darkgray] at (9.55,-3.65) {Pourquoi};
  \node[anchor=north west,text width=5.4cm,align=left] at (9.55,-4.0)
    {Intention : crédit immobilier\\
     Mots-clés : acheter, appartement\\
     Taux d'effort $\leq$ 35\,\% : \textcolor{kored}{\textbf{42,1\,\%}}\\
     Durée $\leq$ 25 ans\\
     Taux indicatif : 3,5\,\% (source, date)\\
     Proposé : rendez-vous conseiller\\
     Écarté : prêt sur 30 ans};
  % Frugalité
  \draw[darkgray!60] (9.45,-7.1) rectangle (15.15,-9.15);
  \node[anchor=north west,font=\scriptsize\bfseries,text=darkgray] at (9.55,-7.15) {Frugalité};
  \node[anchor=north west,text width=5.4cm,align=left] at (9.55,-7.5)
    {Comprendre : mots-clés, \dots{} ms\\
     Conseiller : règles, \dots{} ms\\
     Formuler : \texttt{ministral-3:3b}, \dots{} s\\
     Total : \dots{} s ; \dots{} Wh ; 3,5 Go};
\end{tikzpicture}
\end{adjustbox}
\caption{Maquette de l'écran. Les valeurs sont des exemples ; les points de suspension sont
remplis en direct par les mesures de la trace.}
\label{fig:maquette}
\end{figure}

\paragraph{Accessibilité.} Taille de texte réglable, contraste élevé, mode \og mots simples \fg,
termes du glossaire explicables d'un clic. C'est modeste, mais c'est ce que La Banque Postale a
demandé de ne jamais oublier, et ça se montre en dix secondes.

\paragraph{Les messages de démo.} Une liste latérale contient les messages exacts du script : un
clic les envoie. Pas de faute de frappe devant le jury, pas de temps perdu à taper.

\section{Le script de la démonstration}

Six minutes en direct, sur la machine fournie. Une personne pilote, une autre commente.

\begin{center}\small
\begin{tabularx}{\textwidth}{@{}l L{3.1cm} Y L{2.4cm}@{}}
\toprule
\textbf{Temps} & \textbf{Scène} & \textbf{Ce que le jury voit} & \textbf{Critère} \\
\midrule
0:00 & Inès : \og J'ai perdu ma carte hier soir \fg & L'assistant demande confirmation, met la carte
  en opposition, donne une référence et rassure. Panneau frugalité : \textbf{compris par les
  mots-clés, aucun LLM appelé}. &
  Métier, frugalité \\
1:00 & Léa : l'appartement & Il ne redemande pas les revenus ; il calcule, refuse avec le calcul,
  propose des leviers. Panneau pourquoi : la règle des 35\,\%, la ligne écartée. & Explicabilité \\
2:30 & Sophie puis Lucas : \og Je voudrais mettre de côté \fg & La même phrase, deux clients, deux
  conseils différents, chacun avec ses raisons. & Personnalisation \\
3:30 & Marc : les 8\eur{} de frais & Il retrouve les frais dans l'historique, les explique, cite
  le plafond légal, oriente vers l'offre clientèle fragile. & Métier, inclusion \\
4:30 & Julie : \og Votre appli ne marche jamais ! \fg & Il reconnaît la frustration, ne se
  justifie pas, transmet au conseiller un résumé des trois contacts. & Métier \\
5:15 & Le bilan & Panneau frugalité cumulé ; graphique des configurations A, B et C. & Frugalité \\
\bottomrule
\end{tabularx}
\end{center}

Paul et le virement en mots simples sont dans la vidéo de secours et servent de réponse si le jury
demande \og et pour un client âgé ? \fg.

\section{Le plan de secours}

\begin{itemize}
  \item \textbf{Une vidéo} de tout le script, enregistrée mercredi matin sur la machine de démo.
  \item \textbf{Un mode gabarits} : un interrupteur coupe le petit LLM. La démo continue, moins
        fluide mais exacte. C'est aussi la configuration C, donc ce n'est pas un aveu d'échec.
  \item \textbf{Une vérification réelle} d'Ollama au lancement de l'interface (pas seulement la
        présence d'une configuration) et des modèles préchargés avant l'entrée du jury.
  \item \textbf{Une répétition sur la machine de démo} mercredi entre 12h et 13h15.
\end{itemize}
